\documentclass[10pt,a4paper]{amsart}
\usepackage[margin=19mm]{geometry}
\usepackage{amsmath,amssymb,mathtools}
\usepackage[scr=rsfs,cal=euler]{mathalfa}
\usepackage{xfrac}
\usepackage[unicode,hidelinks,bookmarks=false]{hyperref}
\newcommand{\ie}{i.e.\ }
\newcommand{\cf}{cf.\ }
\newcommand{\resp}{respectively\ }
\newcommand{\Subsection}{\subsection}
\providecommand{\nobreakdash}{\nobreak\hbox{-}}
\setlength{\emergencystretch}{2em}
\multlinegap0pt
\usepackage{amssymb}
\usepackage{mathtools}
\usepackage{enumerate}
\usepackage{dsfont}
\usepackage{xcolor}
\newtheorem{lemma}{Lemma}
\newcommand{\N}{\mathbb N}
\title{On the plasticity of the unit spheres of $\ell_1$, $\ell_{\infty}$, $c$, and Hilbert spaces}
\author{Maksym Levchenko, Olesia Zavarzina}
\date{Source: arXiv:2605.06622v1 (CC BY 4.0)}
\begin{document}
\maketitle

\noindent\emph{Source and licence.} On the plasticity of the unit spheres of $\ell_1$, $\ell_{\infty}$, $c$, and Hilbert spaces. Version arXiv:2605.06622v1 (CC BY 4.0), distributed by arXiv under the Creative Commons Attribution 4.0 International licence, http://creativecommons.org/licenses/by/4.0/, as recorded in the arXiv licence metadata for this exact version and shown on the abstract page https://arxiv.org/abs/2605.06622v1. Full licence notice: \texttt{licenses/arxiv-2605.06622-CC-BY-4.0.txt}.\par\smallskip

\noindent\textit{Editorial scope.} Counted unit: the article's lemma lem\_main (source0.tex lines 367--369). The article's complete original proof of it is delivered with it (source0.tex lines 370--421, one \texttt{proof} environment of 52 lines). No mathematics is written, repaired or re-derived: the statement and the proof are the article's own lines, in the article's own notation, and no retained line is substituted or re-flowed.  No further source-local context is needed: every cross-reference of every retained line resolves inside the two delivered spans.  Nothing was omitted from any retained span, so the package carries no omission record.  No retained line contains a citation, so the wrapper carries no bibliography and no bibliography entry.  No retained line contains a figure environment, an image-inclusion command or drawing code, so no image is delivered and the package declares no asset dependency.  Editorial formatting only, in addition: an AMS \texttt{amsart} preamble that restates the article's own declarations and macros used by the retained lines, namely the environments \texttt{lemma} on separate counters, together with the standard packages those lines need.  The retained environments print in this document as Lemma 1 in order of appearance, whereas the article prints the same items with the numbering of its own full document.  The complete native source of the article as distributed in the arXiv e-print for this exact version is bundled unchanged as \texttt{sources/200039/source0.tex}.
\begin{lemma}\label{lem_main}
    Let $F: S_X\rightarrow S_X$ be a non-expansive bijection. Then the elements of $M$ are fixed points of $I_F$.
\end{lemma}

\begin{proof}
    We are going to prove by induction that the elements of $M_n$ are fixed points of $I_F$ for every natural $n$.
    Let $x \in M_1$, $y = I_F(x)$. Then there is an $r \in (-1,1)$ such that $x^k=r$ for every $k\in \N \setminus \{N\}$ for some $N \in \N$. If $r=0$, then $x=\pm e_N$, and thus $x$ is a fixed point of $I_F$, so we may WLOG assume that $|r|\in (0,1)$. It is implied by (5) that $|y^k|\leq |r|<1$ for every $k \in \N \setminus \{N\}$, so $|y^N|=1=|x^N|$, and (6) means that $y^N=x^N$. Now let $p\in \mathtt{ext}(B_X)$, $p^k=\mathtt{sgn}(x^k)$. The non-expansiveness of $I_F$ and (6) together mean that $||p-x||=1-|r|\geq ||p-y||\geq 1-|y^k|$ for every $k \in \N$, thus $x=y=I_F(x)$.
    Now assume that for some natural $n \geq 2$ all the elements of $M_{n-1}$ are fixed points of $I_F$, and let $x \in M_{n}$, $y=I_F(x)$. Let $A_0,A_1,...A_n\subset \N$ and $r_1,r_2,...,r_n\in [0, 1)$ be such that $|x^i|=r_k$ for all $i \in A_k$, $1 \leq k \leq n$, with conditions (7-9) satisfied. First we are going to demonstrate that $x^i=y^i$ for all $i\in A_k$, $2 \leq k \leq n$. Let $2 \leq k \leq n$, and define $\tilde{x} \in S_X$ by 
     \[\tilde{x}^i=\begin{cases} \mathtt{sgn}(x^i) \frac12  (r_{k-1}+r_k), \quad i \in A_{k-1} \sqcup A_k,\\
    x^i, \quad i \in \N \setminus (A_{k-1} \sqcup A_k).
    \end{cases}\]
    Then $\tilde{x}\in M_{n-1}$, so $\tilde{x}$ is a fixed point of $I_F$, and \[||\tilde{x}-y|| \leq ||\tilde{x}-x||=|\tilde{x}^i|-|x^i|\leq |\tilde{x}^i|-|y^i| \leq ||\tilde{x}-y||\]
    for all $i \in A_k$.\\
    Now all that remains is to show that $x^i=y^i$ for every $i \in A_1$. 
    Suppose in contradiction that $|x^i|>|y^i|$ for some $i \in A_1$, and let $\lambda \in (r_1,1)$ be such that $1-\lambda <|x^i|-|y^i|$. Now let $p \in M_n$ be defined by
    \[p^j=\begin{cases} \lambda \mathtt{sgn}(x^j) , \quad j \in A_1,\\
    x^j, \quad j \in \N \setminus A_1,
    \end{cases}\]
    and let $q=I_F^{-1}(p)$.
    We are going to show that
   \begin{equation}\label{eqns}
       q^j=p^j=x^j=y^j
   \end{equation} 
    for every $j \in \N \setminus A_1$, so that 
    \begin{equation*}
        ||p-q||=\sup_{j\in A_1} |p^j-q^j|\leq 1-\lambda.
    \end{equation*}
    Fix some arbitrary $j \in \N \setminus (A_1 \sqcup A_0)$, and let $\varepsilon \in (0,1-\lambda)$. Now define $\tilde p \in M_{n-1}$ by 
    \[\tilde{p}^l=\begin{cases} \mathtt{sgn}(p^l) , \quad l=j,\\
    (\lambda+\varepsilon) \mathtt {sgn}(x^l), \quad l \in A_0 \sqcup A_1 \sqcup (A_2 \setminus  \{j\}),\\
    p^l + \varepsilon \mathtt{sgn}(p^l), \quad l \in A_k \setminus \{j\}, \ k \geq 3.
    \end{cases}\]
Then 
\[|p^l-\tilde{p} ^l|=\begin{cases} 1-|p^l| , \quad l=j,\\
   1- \lambda-\varepsilon , \quad l \in A_0 ,\\
    \lambda +\varepsilon -|p^l|, \quad l \in A_2 \setminus \{j\}, \\
    \varepsilon, \quad l \in A_1 \sqcup(A_k \setminus \{j\}), \ k \geq 3,
    \end{cases}\]
so $|p^j-\tilde{p}^j|=1-|p^j|>|p^l-\tilde{p} ^l|$ for every $l \neq j$. If $j \in A_k$, $2 \leq k \leq n$, and $l \in A_m$, $0 \leq m \leq k$, $l \neq j$, then 
\begin{equation*}
    |q^l-\tilde p ^l| \leq \begin{cases}
        1 - \lambda - \varepsilon, \quad l \in A_0,\\
        \max (\varepsilon, 1-\lambda-\varepsilon), \quad l \in A_1,\\
        \max(\lambda +\varepsilon -r_2, 1-\lambda-\varepsilon), \quad l \in A_2,\\
        \max(\varepsilon,1-r_m-\varepsilon), \quad l \in A_m, \ 2<m\leq k,
    \end{cases}
\end{equation*}
so
\begin{equation}
    |q^l-\tilde p^l|<1-|p_j|=|p^l-\tilde p^j|=||p-\tilde p||.
\end{equation}
Thus (10) follows from (11) for $j \in A_n$, and for $j \in A_k$, $2 \leq k <n$, it follows from (11) and the fact that $q^{l}=p^{l}$ for every $l \in \bigsqcup_{m=k+1}^nA_m$.\\
Now we can finally complete the proof by contradiction:
\[||q-x||\leq||q-p||+||p-x||\leq1-\lambda+||p-x||<|x^i|-|y^i|+||p-x|| \leq||p-y||=||I_F(q)-I_F(x)||.\]
   The strict inequality in the chain above contradicts the fact that $I_F$ is non-expansive.
\end{proof}

\end{document}
